% -*- coding: utf-8 -*-
\chapter{Conclusiones}

El presente Trabajo de Fin de Máster se planteó con la finalidad de evaluar, medir y validar cuantitativamente la capacidad de la arquitectura Retrieval-Augmented Generation (RAG) para mitigar el problema de las alucinaciones factuales en Modelos de Lenguaje de Gran Escala (LLMs) aplicados a la documentación normativa y técnica del sector agrícola, centrándose específicamente en la Política Agraria Común (PAC) y el Programa POSEI. Tras la ejecución del protocolo experimental, se ha alcanzado la consecución plena y sistemática de la totalidad de las metas trazadas en la fase de planificación del proyecto.

En relación con el primer objetivo específico, centrado en la delimitación, ingestión y estructuración del corpus documental normativo agrícola, este se ha cumplido al 100\%. Se logró consolidar un corpus no paramétrico representativo compuesto por \$N = 20\$ documentos oficiales de la PAC y el POSEI, con una extensión equivalente a \$P = 50\$ páginas de alta densidad técnica e iconográfica. La integración de herramientas de análisis de maquetación avanzada con soporte OCR permitió preservar la jerarquía de títulos, artículos, anexos y tablas de coeficientes de admisibilidad en formato estructurado, superando de manera efectiva las limitaciones inherentes al procesamiento de textos mediante extracción plana de archivos PDF.

Respecto al segundo objetivo específico, orientado al diseño e implementación de una arquitectura RAG con anclaje estricto y trazabilidad por metadatos, se alcanzó igualmente un nivel de cumplimiento del 100\%. Para ello, se desplegó una tubería de recuperación híbrida que combina la búsqueda por similitud vectorial densa con la búsqueda léxica \$BM25\$ y un módulo de reordenamiento (re-ranking) basado en Cross-Encoders. Asimismo, el enriquecimiento de cada fragmento mediante metadatos granulares que identifican el documento, la sección y la página de origen permitió garantizar la citación exacta de las fuentes documentales en cada una de las respuestas generadas por el modelo.

El tercer objetivo específico, referente a la construcción de un conjunto de datos de prueba acotado y balanceado en el dominio agrario, se cumplió al 100\% mediante la elaboración y validación de un dataset sintético compuesto por 150 pruebas. Estas pruebas se distribuyeron homogéneamente en siete categorías funcionales de complejidad creciente, abarcando desde factores directos y consultas temporales hasta análisis comparativos, cálculos numéricos, relaciones inter-normativas, amplitud de contexto y análisis holísticos, proporcionando así una base empírica rigurosa para la auditoría del sistema.

En cuanto al cuarto objetivo específico, correspondiente a la evaluación cuantitativa desacoplada entre el componente recuperador y el generador, se logró un cumplimiento del 100\% gracias a la aplicación de las métricas formales del marco RAGAS. La medición de parámetros como Fidelidad, Relevancia de la Respuesta, Precisión del Contexto, Exhaustividad del Contexto y Calidad de Citación proporcionó indicadores objetivos que permitieron aislar los errores propios del módulo de recuperación de las alucinaciones generadas por el modelo de lenguaje.

Finalmente, la consecución de los cuatro objetivos específicos permitió alcanzar al 100\% el objetivo general del proyecto, demostrando empíricamente en qué medida la arquitectura RAG reduce la tasa de alucinaciones y mejora la exactitud factual frente a un LLM standalone. El sistema AgroRAG alcanzó una puntuación media de Exactitud Factual de , una Relevancia de  y un rendimiento de recuperación con un  global de . Estos hallazgos confirman que, mientras los modelos aislados sufren de alucinaciones extrínsecas debido a las restricciones de su memoria paramétrica, la arquitectura RAG condiciona con éxito la generación a la evidencia documental recuperada, reduciendo drásticamente los errores en la gestión de normativa agrícola.

\section{Limitaciones}
A pesar de haber alcanzado el total de los objetivos propuestos, la ejecución experimental permitió identificar diversos cuellos de botella y limitaciones metodológicas que impidieron que el sistema alcanzara un rendimiento perfecto. La principal limitación técnica del proyecto radica en la degradación observada ante consultas de gran amplitud normativa y de síntesis holística. La estrategia de fragmentación estática tradicional provocó que, en consultas cuya respuesta requería interconectar artículos dispersos a lo largo de varios capítulos o anexos del texto legal, el recuperador denso registrara una caída pronunciada en su rendimiento, con un \$MRR\$ de \$0\{,\}4700\$. Este desacoplamiento en la recuperación entregó información parcial al generador, derivando en respuestas compuestas por abstenciones excesivas o en la emisión de fragmentos incompletos.

Asimismo, se detectó una inestabilidad latente en el procesamiento de tablas normativas complejas y estructuras multidimensionales de cálculo numérico. A pesar del empleo de motores de reconocimiento óptico de caracteres, el recuperador vectorial fragmentó de manera inconsistente las celdas combinadas de los escalados de ayuda del POSEI, perdiendo la relación de dependencia entre filas y columnas y afectando la exhaustividad del contexto. Por último, la auditoría automática evidenció una limitación metodológica en el uso de LLMs como jueces evaluadores, observándose un sesgo sistemático que favorecía respuestas de gran extensión y redacción burocrática sobre aquellas que ofrecían precisión matemática estricta.

\section{Prospectiva}
Con el propósito de dar continuidad a esta línea de investigación y superar las barreras técnicas identificadas, se plantea como principal vía de desarrollo futuro la transición desde el RAG vectorial plano hacia arquitecturas basadas en Grafos de Conocimiento (GraphRAG). Modelar la normativa agrícola como una red explícita de entidades legales, artículos, jerarquías y excepciones permitirá ejecutar búsquedas relacionales multi-salto, resolviendo de manera estructural la dispersión contextual observada en las consultas transversales.

Adicionalmente, se contempla la evolución del sistema hacia un esquema de RAG Agéntico (Agentic RAG), en el cual diversos agentes especializados en ruteo, cálculo numérico y verificación jurídica colaboren de forma orquestada para validar la consistencia lógica y matemática de las respuestas antes de emitir un dictamen final. Finalmente, se propone escalar la base de conocimiento no paramétrica para integrar la totalidad de los boletines oficiales autonómicos y reglamentos de la Unión Europea.

\section{Consideraciones finales}
Los hallazgos derivados de este trabajo demuestran que las arquitecturas RAG no constituyen una solución mágica que elimina por completo el riesgo de alucinación, sino un marco de control probabilístico de alta efectividad. En dominios de elevada responsabilidad como la gestión de ayudas públicas de la PAC y el POSEI, la seguridad jurídica exige trascender la generación del lenguaje para consolidar sistemas transparentes, auditables y fundamentados en metadatos granulares.

A nivel personal y profesional, la elaboración de esta investigación ha consolidado mi dominio de las arquitecturas avanzadas de Procesamiento del Lenguaje Natural y el aprendizaje profundo. Más allá del desarrollo técnico, el mayor aprendizaje ha residido en la asimilación del rigor del método científico, comprendiendo la necesidad de auditar críticamente las métricas globales para identificar las vulnerabilidades subyacentes en los sistemas de inteligencia artificial.

El recorrido a lo largo del Máster ha significado una transformación integral en mi perfil técnico, permitiéndome conectar la ingeniería de datos con la gobernanza tecnológica y la resolución de problemas reales. Esta formación me ha dotado de las herramientas necesarias para diseñar soluciones innovadoras que aporten valor tangible a la sociedad y al sector productivo.

La experiencia acumulada en la estructuración de este proyecto ha fortalecido mi percepción sobre la importancia de la función docente e investigadora. La capacidad de simplificar y transmitir conceptos matemáticos y computacionales de gran complejidad resulta indispensable en el entorno académico contemporáneo. En este sentido, considero prioritario impulsar metodologías de enseñanza centradas en la auditoría técnica y la evaluación crítica de los modelos emergentes.

Expreso mi sincero agradecimiento al equipo tutelar por su valioso acompañamiento, paciencia y rigor académico durante el desarrollo de esta tesis. De igual forma, extiendo mi gratitud al cuerpo docente del Máster por la excelencia de la enseñanza impartida, así como a mi familia y compañeros por su apoyo constante e incondicional en cada etapa de este proceso.

No puedo finalizar este trabajo sin manifestar mi ferviente deseo que la transformación digital de nuestro sector agrario avance decididamente hacia un modelo transparente, justo y accesible. La simplificación de las cargas burocráticas mediante tecnologías seguras constituye una deuda histórica con nuestros agricultores y ganaderos, y confío en que la aportación de este Trabajo de Fin de Máster sea un paso firme hacia la consolidación de una Inteligencia Artificial al servicio del bienestar de nuestro campo.



